\section{引言}

作为强相互作用的规范理论，QCD 的一个核心挑战是直接预言强子结构。特别是，部分子分布函数（PDF）与广义部分子分布（GPD）直到最近仍是第一性原理计算难以处理的问题。尽管 PDF 与 GPD 的完整 \textit{ab initio}（从头）确定尚未出现，文献~\cite{Ji:2013dva} 最近提出了一种有前景的新方法。在该方法中，PDF 与 GPD 从有限动量核子态之间空间延展算符的矩阵元中提取。这些矩阵元通常称为准部分子分布（quasi-PDF）或准分布。相关框架也已被提出：文献~\cite{Ma:2014jla,Ma:2014jga,Ma:2017pxb} 把准分布视为可从中因子化出光前 PDF 的``格点截面''，而文献~\cite{Radyushkin:2016hsy,Radyushkin:2017cyf,Orginos:2017kos} 引入并研究了密切相关的赝分布。

准分布是时间定域算符的矩阵元，因此可以直接用格点 QCD 计算 \cite{Carlson:2017gpk,Briceno:2017cpo}，其中 QCD 离散化在欧氏超立方格点上。准分布与赝分布的初步非微扰结果令人鼓舞 \cite{Lin:2014zya,Alexandrou:2015rja,Chen:2016utp,Orginos:2017kos,Zhang:2017bzy,Alexandrou:2017qpu}。此外，许多理论问题 \cite{Li:2016amo,Rossi:2017muf,Carlson:2017gpk} 已被澄清或解决，包括空间 Wilson 线算符可乘重正化的证明 \cite{Ji:2017oey,Ishikawa:2017faj}、光前 PDF 从准分布因子化的证明 \cite{Ma:2014jla,Ma:2014jga}，以及从欧氏关联函数提取的矩阵元与通过闵氏空间 Lehmann-Symanzik-Zimmermann 约化程序定义的矩阵元相同的证明 \cite{Briceno:2017cpo,Xiong:2017jtn,Ji:2017rah}。

必须解决的计算挑战之一是有限格点调节子在准分布与赝分布中诱导的幂次发散。定义这些分布的 Wilson 线算符具有随 Wilson 线长度除以格距而指数增长的发散，该发散必须非微扰地移除。已有若干方案被提出：文献~\cite{Chen:2016fxx} 与 \cite{Ishikawa:2016znu} 的作者建议通过指数化的质量重正化移除幂次发散；最近，正则化不变动量减除方案（RI/MOM \cite{Chen:2017mzz} 与 RI' \cite{Alexandrou:2017huk}）被建议作为非微扰重正化程序。在文献~\cite{Monahan:2016bvm} 中，我们引入了另一种方法——涂抹准分布，它利用梯度流的性质 \cite{Narayanan:2006rf,Luscher:2011bx,Luscher:2013cpa} 规避了幂次发散的部分挑战。

梯度流是原始夸克与胶子自由度在一个新参数——流时——中的经典演化，或称单参数映射。流时指数式地压制紫外场涨落，这对应于在实空间中涂抹原始自由度。梯度流的关键性质是：除一个可乘的费米子波函数重正化 \cite{Luscher:2013cpa} 之外，原理论的有限关联函数在非零流时保持有限 \cite{Luscher:2011bx}。把流时固定在物理单位，就可以保证在有限格距下非微扰确定的矩阵元在连续极限下保持有限。因此，梯度流提供了一种非微扰的、规范不变的方法，使准分布即使在连续极限下也有限 \cite{Monahan:2016bvm}。所得的连续矩阵元随后可直接与光前 PDF 联系，或与例如在 $\ms$ 方案下重正化的准分布或赝分布联系。实践中，把这个匹配程序分几步进行更为简单：先把涂抹分布与未涂抹分布联系起来，再利用未涂抹分布与光前 PDF 之间的已知关系 \cite{Xiong:2013bka,Ji:2015jwa,Wang:2017qyg,Stewart:2017tvs}。$\ms$ 方案下相关矩阵元的微扰计算见文献~\cite{Constantinou:2017sej}，而本文的重点是确定有限流时下的矩阵元。

我在微扰论单圈水平研究定义涂抹准分布的矩阵元。在静止的外部无质量夸克态之间计算涂抹 Wilson 线算符的矩阵元，使得所有积分都能解析求值。单圈水平下，Wilson 线幂次发散表现为在 $\oz \gg 1$ 时对 $\oz = z/r_\tau$ 线性的贡献，其中 $z$ 是 Wilson 线长度，$r_\tau$ 是梯度流涂抹半径。减去这一贡献后，其余矩阵元有限，且具有良好定义的 $\oz\to0$ 极限。在小流时区域（即 $\oz \gg 1$），矩阵元仅对 $\oz$ 对数依赖，并满足与通常重正化群方程类似的关系。

第~\ref{sec:defns} 节先简要回顾光前 PDF、涂抹准分布及其相互关系。第~\ref{sec:pertth} 节讨论相关矩阵元的微扰计算，第~\ref{sec:numres} 节给出若干数值结果，第~\ref{sec:summary} 节作总结。
